% ══════════════════════════════════════════════════════════════════════════════
%  المشترك — التمهيد الموحّد لوثائق الدروس (مذكرة / وثيقة التلميذ / أعمال تطبيقية)
%  يُستدعى بـ: % ══════════════════════════════════════════════════════════════════════════════
%  المشترك — التمهيد الموحّد لوثائق الدروس (مذكرة / وثيقة التلميذ / أعمال تطبيقية)
%  يُستدعى بـ: % ══════════════════════════════════════════════════════════════════════════════
%  المشترك — التمهيد الموحّد لوثائق الدروس (مذكرة / وثيقة التلميذ / أعمال تطبيقية)
%  يُستدعى بـ: % ══════════════════════════════════════════════════════════════════════════════
%  المشترك — التمهيد الموحّد لوثائق الدروس (مذكرة / وثيقة التلميذ / أعمال تطبيقية)
%  يُستدعى بـ: \input{../_common/preamble.tex}
%  يُصرَّف بـ XeLaTeX حصراً (عربية + RTL).
% ══════════════════════════════════════════════════════════════════════════════

% ── Geometry ─────────────────────────────────────────────────────────────────
\usepackage[
  top=3.2cm, bottom=2.5cm, left=2.2cm, right=2.2cm,
  headheight=26pt
]{geometry}

% ── Packages: load fancyhdr and friends BEFORE polyglossia ───────────────────
\usepackage{fancyhdr}
\usepackage{titlesec}
\usepackage{enumitem}
\usepackage{xcolor}
\usepackage{tikz}
\usepackage{eso-pic}
\usepackage{tcolorbox}
\tcbuselibrary{skins, breakable}
\usepackage{array}
\usepackage{longtable}
\usepackage{colortbl}
\usepackage{multicol}
\usepackage{multirow}
\usepackage{hyperref}

% ── Arabic / RTL — load AFTER fancyhdr ───────────────────────────────────────
\usepackage{polyglossia}
\setmainlanguage[numerals=western]{arabic}
\setotherlanguage{english}

% ── Fonts ────────────────────────────────────────────────────────────────────
\setmainfont{Amiri}[
  BoldFont = * Bold,
  ItalicFont = * Italic,
]
\newfontfamily\arabicfont{Amiri}[
  Script = Arabic,
  BoldFont = * Bold,
]
\setmonofont{DejaVu Sans Mono}
\newfontfamily\arabicfonttt{Amiri}[Script = Arabic, BoldFont = * Bold]

% ── Colors (project palette) ─────────────────────────────────────────────────
\definecolor{primary}{HTML}{1A5276}
\definecolor{secondary}{HTML}{2E86C1}
\definecolor{lightbg}{HTML}{EBF5FB}
\definecolor{darktext}{HTML}{1C2833}
\definecolor{rulegray}{HTML}{ABB2B9}
\definecolor{success}{HTML}{1E8449}

% ── Header / Footer ──────────────────────────────────────────────────────────
\pagestyle{fancy}
\fancyhf{}
\renewcommand{\headrulewidth}{0.8pt}
\renewcommand{\footrulewidth}{0.4pt}
\fancyhead[L]{\small\color{primary} المعلوماتية — السنة الأولى ثانوي}
\fancyhead[R]{\small\color{primary} ثانوية الشهيد حكومي العيد أدرار}
\fancyfoot[C]{\small\color{rulegray}\thepage}
\fancyfoot[R]{\footnotesize\color{rulegray} 2026/2027}
\fancyfoot[L]{\footnotesize\color{rulegray} الأستاذ: لفقير عبدالجليل}

% ── Page Border: 1cm from all four edges ─────────────────────────────────────
\AddToShipoutPictureBG{%
  \begin{tikzpicture}[remember picture, overlay]
    \draw[primary, line width=1.5pt]
      ([xshift=-1cm, yshift=-1cm]current page.north east)
      rectangle
      ([xshift=1cm, yshift=1cm]current page.south west);
  \end{tikzpicture}%
}

% ── Section formatting ───────────────────────────────────────────────────────
\titleformat{\section}
  {\large\bfseries\color{primary}}
  {\colorbox{primary}{\color{white}\thesection}}
  {0.7em}{}
  [\vspace{-0.15em}{\color{primary}\titlerule[1.2pt]}]

\titleformat{\subsection}
  {\normalsize\bfseries\color{secondary}}{\thesubsection}{0.5em}{}

% عنوان قسم بدون ترقيم (لا يُرسم له مربّع رقم فارغ)
\newcommand{\plainsec}[1]{%
  \par\vspace{6pt}%
  {\large\bfseries\color{primary} #1}\par
  \vspace{-0.35em}{\color{primary}\hrule height 1.2pt}\par\vspace{3pt}%
}

% ── Reusable boxes ───────────────────────────────────────────────────────────
% صندوق معلومات عام
\newtcolorbox{infobox}[1][]{
  enhanced, colback=lightbg, colframe=secondary,
  coltitle=white, fonttitle=\bfseries, title=#1,
  boxrule=0.9pt, arc=4pt, left=10pt, right=10pt, top=8pt, bottom=8pt,
  breakable,
}
% صندوق "أتذكّر" — الخلاصة
\newtcolorbox{rememberbox}[1][أتذكّر]{
  enhanced, colback=lightbg, colframe=primary,
  coltitle=white, fonttitle=\bfseries, title=#1,
  boxrule=1pt, arc=4pt, left=10pt, right=10pt, top=8pt, bottom=8pt,
  breakable,
}
% صندوق نشاط
\newtcolorbox{activitybox}[1][نشاط]{
  enhanced, colback=white, colframe=secondary,
  coltitle=white, fonttitle=\bfseries, title=#1,
  boxrule=1pt, arc=4pt, left=10pt, right=10pt, top=8pt, bottom=8pt,
  breakable,
}
% صندوق الإشكالية
\newtcolorbox{problembox}[1][الإشكالية]{
  enhanced, colback=white, colframe=primary,
  coltitle=white, fonttitle=\bfseries, title=#1,
  boxrule=1pt, arc=4pt, left=10pt, right=10pt, top=8pt, bottom=8pt,
  breakable,
}
% صندوق الواجب المنزلي
\newtcolorbox{homeworkbox}[1][واجب منزلي]{
  enhanced, colback=lightbg, colframe=success,
  coltitle=white, fonttitle=\bfseries, title=#1,
  boxrule=1pt, arc=4pt, left=10pt, right=10pt, top=8pt, bottom=8pt,
  breakable,
}

% ── Lists ────────────────────────────────────────────────────────────────────
\newlist{numlist}{enumerate}{1}
\setlist[numlist,1]{
  label=\arabic*.\hspace{0.8em},
  labelwidth=1.4em, left=0em .. 1.4em,
  itemsep=4pt, parsep=0pt, itemindent=0pt,
}
\newlist{bullets}{itemize}{1}
\setlist[bullets,1]{
  label=\textcolor{secondary}{\textbullet},
  left=0.2em .. 1.2em,
  itemsep=3pt, parsep=0pt,
}

% ── Helpers ──────────────────────────────────────────────────────────────────
% سطر نقاط للإجابة (يجب أن يكون داخل صندوق ليرتسم)
\newcommand{\answerline}{\par\vspace{1pt}\noindent\makebox[\linewidth]{\dotfill}}
% عدّة أسطر إجابة
\newcommand{\answerlines}[1]{\begingroup\count255=0 \loop\ifnum\count255<#1 \answerline\advance\count255 by 1 \repeat\endgroup}
% فراغ داخل السطر
\newcommand{\blank}[1][2.2cm]{\underline{\hspace{#1}}}
% مربّع اختيار
\newcommand{\cb}{\raisebox{0.15ex}{\fbox{\rule{0pt}{1.1ex}\hspace{1em}}}}

% ── Lesson title block ───────────────────────────────────────────────────────
% \lessonheader{نوع الوثيقة}{عنوان الدرس}{المجال التعلمي}{الوحدة التعليمية}
\newcommand{\lessonheader}[4]{%
  \begin{center}
    {\small\color{secondary} الجمهورية الجزائرية الديمقراطية الشعبية}\\[1pt]
    {\small\color{secondary} وزارة التربية الوطنية \quad$\cdot$\quad مديرية التربية لولاية أدرار \quad$\cdot$\quad ثانوية الشهيد حكومي العيد أدرار}\\[5pt]
    {\Large\bfseries\color{primary} #1}\\[3pt]
    {\large\bfseries\color{primary} #2}
  \end{center}
  \vspace{-2pt}
  {\color{secondary}\hrule height 1.2pt}
  \vspace{4pt}
  \begin{center}
  \begin{tabular}{@{}>{\bfseries\color{primary}}r@{\hspace{0.5em}}l@{\hspace{2em}}>{\bfseries\color{primary}}r@{\hspace{0.5em}}l@{}}
    المجال التعلمي : & #3 & الوحدة التعليمية : & #4 \\
    المستوى : & السنة الأولى ثانوي & المدة : & ساعة واحدة \\
  \end{tabular}
  \end{center}
  \vspace{2pt}
  {\color{secondary}\hrule height 1.2pt}
  \vspace{6pt}
}

\hypersetup{
  colorlinks=false,
  pdfauthor={لفقير عبدالجليل},
}

%  يُصرَّف بـ XeLaTeX حصراً (عربية + RTL).
% ══════════════════════════════════════════════════════════════════════════════

% ── Geometry ─────────────────────────────────────────────────────────────────
\usepackage[
  top=3.2cm, bottom=2.5cm, left=2.2cm, right=2.2cm,
  headheight=26pt
]{geometry}

% ── Packages: load fancyhdr and friends BEFORE polyglossia ───────────────────
\usepackage{fancyhdr}
\usepackage{titlesec}
\usepackage{enumitem}
\usepackage{xcolor}
\usepackage{tikz}
\usepackage{eso-pic}
\usepackage{tcolorbox}
\tcbuselibrary{skins, breakable}
\usepackage{array}
\usepackage{longtable}
\usepackage{colortbl}
\usepackage{multicol}
\usepackage{multirow}
\usepackage{hyperref}

% ── Arabic / RTL — load AFTER fancyhdr ───────────────────────────────────────
\usepackage{polyglossia}
\setmainlanguage[numerals=western]{arabic}
\setotherlanguage{english}

% ── Fonts ────────────────────────────────────────────────────────────────────
\setmainfont{Amiri}[
  BoldFont = * Bold,
  ItalicFont = * Italic,
]
\newfontfamily\arabicfont{Amiri}[
  Script = Arabic,
  BoldFont = * Bold,
]
\setmonofont{DejaVu Sans Mono}
\newfontfamily\arabicfonttt{Amiri}[Script = Arabic, BoldFont = * Bold]

% ── Colors (project palette) ─────────────────────────────────────────────────
\definecolor{primary}{HTML}{1A5276}
\definecolor{secondary}{HTML}{2E86C1}
\definecolor{lightbg}{HTML}{EBF5FB}
\definecolor{darktext}{HTML}{1C2833}
\definecolor{rulegray}{HTML}{ABB2B9}
\definecolor{success}{HTML}{1E8449}

% ── Header / Footer ──────────────────────────────────────────────────────────
\pagestyle{fancy}
\fancyhf{}
\renewcommand{\headrulewidth}{0.8pt}
\renewcommand{\footrulewidth}{0.4pt}
\fancyhead[L]{\small\color{primary} المعلوماتية — السنة الأولى ثانوي}
\fancyhead[R]{\small\color{primary} ثانوية الشهيد حكومي العيد أدرار}
\fancyfoot[C]{\small\color{rulegray}\thepage}
\fancyfoot[R]{\footnotesize\color{rulegray} 2026/2027}
\fancyfoot[L]{\footnotesize\color{rulegray} الأستاذ: لفقير عبدالجليل}

% ── Page Border: 1cm from all four edges ─────────────────────────────────────
\AddToShipoutPictureBG{%
  \begin{tikzpicture}[remember picture, overlay]
    \draw[primary, line width=1.5pt]
      ([xshift=-1cm, yshift=-1cm]current page.north east)
      rectangle
      ([xshift=1cm, yshift=1cm]current page.south west);
  \end{tikzpicture}%
}

% ── Section formatting ───────────────────────────────────────────────────────
\titleformat{\section}
  {\large\bfseries\color{primary}}
  {\colorbox{primary}{\color{white}\thesection}}
  {0.7em}{}
  [\vspace{-0.15em}{\color{primary}\titlerule[1.2pt]}]

\titleformat{\subsection}
  {\normalsize\bfseries\color{secondary}}{\thesubsection}{0.5em}{}

% عنوان قسم بدون ترقيم (لا يُرسم له مربّع رقم فارغ)
\newcommand{\plainsec}[1]{%
  \par\vspace{6pt}%
  {\large\bfseries\color{primary} #1}\par
  \vspace{-0.35em}{\color{primary}\hrule height 1.2pt}\par\vspace{3pt}%
}

% ── Reusable boxes ───────────────────────────────────────────────────────────
% صندوق معلومات عام
\newtcolorbox{infobox}[1][]{
  enhanced, colback=lightbg, colframe=secondary,
  coltitle=white, fonttitle=\bfseries, title=#1,
  boxrule=0.9pt, arc=4pt, left=10pt, right=10pt, top=8pt, bottom=8pt,
  breakable,
}
% صندوق "أتذكّر" — الخلاصة
\newtcolorbox{rememberbox}[1][أتذكّر]{
  enhanced, colback=lightbg, colframe=primary,
  coltitle=white, fonttitle=\bfseries, title=#1,
  boxrule=1pt, arc=4pt, left=10pt, right=10pt, top=8pt, bottom=8pt,
  breakable,
}
% صندوق نشاط
\newtcolorbox{activitybox}[1][نشاط]{
  enhanced, colback=white, colframe=secondary,
  coltitle=white, fonttitle=\bfseries, title=#1,
  boxrule=1pt, arc=4pt, left=10pt, right=10pt, top=8pt, bottom=8pt,
  breakable,
}
% صندوق الإشكالية
\newtcolorbox{problembox}[1][الإشكالية]{
  enhanced, colback=white, colframe=primary,
  coltitle=white, fonttitle=\bfseries, title=#1,
  boxrule=1pt, arc=4pt, left=10pt, right=10pt, top=8pt, bottom=8pt,
  breakable,
}
% صندوق الواجب المنزلي
\newtcolorbox{homeworkbox}[1][واجب منزلي]{
  enhanced, colback=lightbg, colframe=success,
  coltitle=white, fonttitle=\bfseries, title=#1,
  boxrule=1pt, arc=4pt, left=10pt, right=10pt, top=8pt, bottom=8pt,
  breakable,
}

% ── Lists ────────────────────────────────────────────────────────────────────
\newlist{numlist}{enumerate}{1}
\setlist[numlist,1]{
  label=\arabic*.\hspace{0.8em},
  labelwidth=1.4em, left=0em .. 1.4em,
  itemsep=4pt, parsep=0pt, itemindent=0pt,
}
\newlist{bullets}{itemize}{1}
\setlist[bullets,1]{
  label=\textcolor{secondary}{\textbullet},
  left=0.2em .. 1.2em,
  itemsep=3pt, parsep=0pt,
}

% ── Helpers ──────────────────────────────────────────────────────────────────
% سطر نقاط للإجابة (يجب أن يكون داخل صندوق ليرتسم)
\newcommand{\answerline}{\par\vspace{1pt}\noindent\makebox[\linewidth]{\dotfill}}
% عدّة أسطر إجابة
\newcommand{\answerlines}[1]{\begingroup\count255=0 \loop\ifnum\count255<#1 \answerline\advance\count255 by 1 \repeat\endgroup}
% فراغ داخل السطر
\newcommand{\blank}[1][2.2cm]{\underline{\hspace{#1}}}
% مربّع اختيار
\newcommand{\cb}{\raisebox{0.15ex}{\fbox{\rule{0pt}{1.1ex}\hspace{1em}}}}

% ── Lesson title block ───────────────────────────────────────────────────────
% \lessonheader{نوع الوثيقة}{عنوان الدرس}{المجال التعلمي}{الوحدة التعليمية}
\newcommand{\lessonheader}[4]{%
  \begin{center}
    {\small\color{secondary} الجمهورية الجزائرية الديمقراطية الشعبية}\\[1pt]
    {\small\color{secondary} وزارة التربية الوطنية \quad$\cdot$\quad مديرية التربية لولاية أدرار \quad$\cdot$\quad ثانوية الشهيد حكومي العيد أدرار}\\[5pt]
    {\Large\bfseries\color{primary} #1}\\[3pt]
    {\large\bfseries\color{primary} #2}
  \end{center}
  \vspace{-2pt}
  {\color{secondary}\hrule height 1.2pt}
  \vspace{4pt}
  \begin{center}
  \begin{tabular}{@{}>{\bfseries\color{primary}}r@{\hspace{0.5em}}l@{\hspace{2em}}>{\bfseries\color{primary}}r@{\hspace{0.5em}}l@{}}
    المجال التعلمي : & #3 & الوحدة التعليمية : & #4 \\
    المستوى : & السنة الأولى ثانوي & المدة : & ساعة واحدة \\
  \end{tabular}
  \end{center}
  \vspace{2pt}
  {\color{secondary}\hrule height 1.2pt}
  \vspace{6pt}
}

\hypersetup{
  colorlinks=false,
  pdfauthor={لفقير عبدالجليل},
}

%  يُصرَّف بـ XeLaTeX حصراً (عربية + RTL).
% ══════════════════════════════════════════════════════════════════════════════

% ── Geometry ─────────────────────────────────────────────────────────────────
\usepackage[
  top=3.2cm, bottom=2.5cm, left=2.2cm, right=2.2cm,
  headheight=26pt
]{geometry}

% ── Packages: load fancyhdr and friends BEFORE polyglossia ───────────────────
\usepackage{fancyhdr}
\usepackage{titlesec}
\usepackage{enumitem}
\usepackage{xcolor}
\usepackage{tikz}
\usepackage{eso-pic}
\usepackage{tcolorbox}
\tcbuselibrary{skins, breakable}
\usepackage{array}
\usepackage{longtable}
\usepackage{colortbl}
\usepackage{multicol}
\usepackage{multirow}
\usepackage{hyperref}

% ── Arabic / RTL — load AFTER fancyhdr ───────────────────────────────────────
\usepackage{polyglossia}
\setmainlanguage[numerals=western]{arabic}
\setotherlanguage{english}

% ── Fonts ────────────────────────────────────────────────────────────────────
\setmainfont{Amiri}[
  BoldFont = * Bold,
  ItalicFont = * Italic,
]
\newfontfamily\arabicfont{Amiri}[
  Script = Arabic,
  BoldFont = * Bold,
]
\setmonofont{DejaVu Sans Mono}
\newfontfamily\arabicfonttt{Amiri}[Script = Arabic, BoldFont = * Bold]

% ── Colors (project palette) ─────────────────────────────────────────────────
\definecolor{primary}{HTML}{1A5276}
\definecolor{secondary}{HTML}{2E86C1}
\definecolor{lightbg}{HTML}{EBF5FB}
\definecolor{darktext}{HTML}{1C2833}
\definecolor{rulegray}{HTML}{ABB2B9}
\definecolor{success}{HTML}{1E8449}

% ── Header / Footer ──────────────────────────────────────────────────────────
\pagestyle{fancy}
\fancyhf{}
\renewcommand{\headrulewidth}{0.8pt}
\renewcommand{\footrulewidth}{0.4pt}
\fancyhead[L]{\small\color{primary} المعلوماتية — السنة الأولى ثانوي}
\fancyhead[R]{\small\color{primary} ثانوية الشهيد حكومي العيد أدرار}
\fancyfoot[C]{\small\color{rulegray}\thepage}
\fancyfoot[R]{\footnotesize\color{rulegray} 2026/2027}
\fancyfoot[L]{\footnotesize\color{rulegray} الأستاذ: لفقير عبدالجليل}

% ── Page Border: 1cm from all four edges ─────────────────────────────────────
\AddToShipoutPictureBG{%
  \begin{tikzpicture}[remember picture, overlay]
    \draw[primary, line width=1.5pt]
      ([xshift=-1cm, yshift=-1cm]current page.north east)
      rectangle
      ([xshift=1cm, yshift=1cm]current page.south west);
  \end{tikzpicture}%
}

% ── Section formatting ───────────────────────────────────────────────────────
\titleformat{\section}
  {\large\bfseries\color{primary}}
  {\colorbox{primary}{\color{white}\thesection}}
  {0.7em}{}
  [\vspace{-0.15em}{\color{primary}\titlerule[1.2pt]}]

\titleformat{\subsection}
  {\normalsize\bfseries\color{secondary}}{\thesubsection}{0.5em}{}

% عنوان قسم بدون ترقيم (لا يُرسم له مربّع رقم فارغ)
\newcommand{\plainsec}[1]{%
  \par\vspace{6pt}%
  {\large\bfseries\color{primary} #1}\par
  \vspace{-0.35em}{\color{primary}\hrule height 1.2pt}\par\vspace{3pt}%
}

% ── Reusable boxes ───────────────────────────────────────────────────────────
% صندوق معلومات عام
\newtcolorbox{infobox}[1][]{
  enhanced, colback=lightbg, colframe=secondary,
  coltitle=white, fonttitle=\bfseries, title=#1,
  boxrule=0.9pt, arc=4pt, left=10pt, right=10pt, top=8pt, bottom=8pt,
  breakable,
}
% صندوق "أتذكّر" — الخلاصة
\newtcolorbox{rememberbox}[1][أتذكّر]{
  enhanced, colback=lightbg, colframe=primary,
  coltitle=white, fonttitle=\bfseries, title=#1,
  boxrule=1pt, arc=4pt, left=10pt, right=10pt, top=8pt, bottom=8pt,
  breakable,
}
% صندوق نشاط
\newtcolorbox{activitybox}[1][نشاط]{
  enhanced, colback=white, colframe=secondary,
  coltitle=white, fonttitle=\bfseries, title=#1,
  boxrule=1pt, arc=4pt, left=10pt, right=10pt, top=8pt, bottom=8pt,
  breakable,
}
% صندوق الإشكالية
\newtcolorbox{problembox}[1][الإشكالية]{
  enhanced, colback=white, colframe=primary,
  coltitle=white, fonttitle=\bfseries, title=#1,
  boxrule=1pt, arc=4pt, left=10pt, right=10pt, top=8pt, bottom=8pt,
  breakable,
}
% صندوق الواجب المنزلي
\newtcolorbox{homeworkbox}[1][واجب منزلي]{
  enhanced, colback=lightbg, colframe=success,
  coltitle=white, fonttitle=\bfseries, title=#1,
  boxrule=1pt, arc=4pt, left=10pt, right=10pt, top=8pt, bottom=8pt,
  breakable,
}

% ── Lists ────────────────────────────────────────────────────────────────────
\newlist{numlist}{enumerate}{1}
\setlist[numlist,1]{
  label=\arabic*.\hspace{0.8em},
  labelwidth=1.4em, left=0em .. 1.4em,
  itemsep=4pt, parsep=0pt, itemindent=0pt,
}
\newlist{bullets}{itemize}{1}
\setlist[bullets,1]{
  label=\textcolor{secondary}{\textbullet},
  left=0.2em .. 1.2em,
  itemsep=3pt, parsep=0pt,
}

% ── Helpers ──────────────────────────────────────────────────────────────────
% سطر نقاط للإجابة (يجب أن يكون داخل صندوق ليرتسم)
\newcommand{\answerline}{\par\vspace{1pt}\noindent\makebox[\linewidth]{\dotfill}}
% عدّة أسطر إجابة
\newcommand{\answerlines}[1]{\begingroup\count255=0 \loop\ifnum\count255<#1 \answerline\advance\count255 by 1 \repeat\endgroup}
% فراغ داخل السطر
\newcommand{\blank}[1][2.2cm]{\underline{\hspace{#1}}}
% مربّع اختيار
\newcommand{\cb}{\raisebox{0.15ex}{\fbox{\rule{0pt}{1.1ex}\hspace{1em}}}}

% ── Lesson title block ───────────────────────────────────────────────────────
% \lessonheader{نوع الوثيقة}{عنوان الدرس}{المجال التعلمي}{الوحدة التعليمية}
\newcommand{\lessonheader}[4]{%
  \begin{center}
    {\small\color{secondary} الجمهورية الجزائرية الديمقراطية الشعبية}\\[1pt]
    {\small\color{secondary} وزارة التربية الوطنية \quad$\cdot$\quad مديرية التربية لولاية أدرار \quad$\cdot$\quad ثانوية الشهيد حكومي العيد أدرار}\\[5pt]
    {\Large\bfseries\color{primary} #1}\\[3pt]
    {\large\bfseries\color{primary} #2}
  \end{center}
  \vspace{-2pt}
  {\color{secondary}\hrule height 1.2pt}
  \vspace{4pt}
  \begin{center}
  \begin{tabular}{@{}>{\bfseries\color{primary}}r@{\hspace{0.5em}}l@{\hspace{2em}}>{\bfseries\color{primary}}r@{\hspace{0.5em}}l@{}}
    المجال التعلمي : & #3 & الوحدة التعليمية : & #4 \\
    المستوى : & السنة الأولى ثانوي & المدة : & ساعة واحدة \\
  \end{tabular}
  \end{center}
  \vspace{2pt}
  {\color{secondary}\hrule height 1.2pt}
  \vspace{6pt}
}

\hypersetup{
  colorlinks=false,
  pdfauthor={لفقير عبدالجليل},
}

%  يُصرَّف بـ XeLaTeX حصراً (عربية + RTL).
% ══════════════════════════════════════════════════════════════════════════════

% ── Geometry ─────────────────────────────────────────────────────────────────
\usepackage[
  top=3.2cm, bottom=2.5cm, left=2.2cm, right=2.2cm,
  headheight=26pt
]{geometry}

% ── Packages: load fancyhdr and friends BEFORE polyglossia ───────────────────
\usepackage{fancyhdr}
\usepackage{titlesec}
\usepackage{enumitem}
\usepackage{xcolor}
\usepackage{tikz}
\usepackage{eso-pic}
\usepackage{tcolorbox}
\tcbuselibrary{skins, breakable}
\usepackage{array}
\usepackage{longtable}
\usepackage{colortbl}
\usepackage{multicol}
\usepackage{multirow}
\usepackage{hyperref}

% ── Arabic / RTL — load AFTER fancyhdr ───────────────────────────────────────
\usepackage{polyglossia}
\setmainlanguage[numerals=western]{arabic}
\setotherlanguage{english}

% ── Fonts ────────────────────────────────────────────────────────────────────
\setmainfont{Amiri}[
  BoldFont = * Bold,
  ItalicFont = * Italic,
]
\newfontfamily\arabicfont{Amiri}[
  Script = Arabic,
  BoldFont = * Bold,
]
\setmonofont{DejaVu Sans Mono}
\newfontfamily\arabicfonttt{Amiri}[Script = Arabic, BoldFont = * Bold]

% ── Colors (project palette) ─────────────────────────────────────────────────
\definecolor{primary}{HTML}{1A5276}
\definecolor{secondary}{HTML}{2E86C1}
\definecolor{lightbg}{HTML}{EBF5FB}
\definecolor{darktext}{HTML}{1C2833}
\definecolor{rulegray}{HTML}{ABB2B9}
\definecolor{success}{HTML}{1E8449}

% ── Header / Footer ──────────────────────────────────────────────────────────
\pagestyle{fancy}
\fancyhf{}
\renewcommand{\headrulewidth}{0.8pt}
\renewcommand{\footrulewidth}{0.4pt}
\fancyhead[L]{\small\color{primary} المعلوماتية — السنة الأولى ثانوي}
\fancyhead[R]{\small\color{primary} ثانوية الشهيد حكومي العيد أدرار}
\fancyfoot[C]{\small\color{rulegray}\thepage}
\fancyfoot[R]{\footnotesize\color{rulegray} 2026/2027}
\fancyfoot[L]{\footnotesize\color{rulegray} الأستاذ: لفقير عبدالجليل}

% ── Page Border: 1cm from all four edges ─────────────────────────────────────
\AddToShipoutPictureBG{%
  \begin{tikzpicture}[remember picture, overlay]
    \draw[primary, line width=1.5pt]
      ([xshift=-1cm, yshift=-1cm]current page.north east)
      rectangle
      ([xshift=1cm, yshift=1cm]current page.south west);
  \end{tikzpicture}%
}

% ── Section formatting ───────────────────────────────────────────────────────
\titleformat{\section}
  {\large\bfseries\color{primary}}
  {\colorbox{primary}{\color{white}\thesection}}
  {0.7em}{}
  [\vspace{-0.15em}{\color{primary}\titlerule[1.2pt]}]

\titleformat{\subsection}
  {\normalsize\bfseries\color{secondary}}{\thesubsection}{0.5em}{}

% عنوان قسم بدون ترقيم (لا يُرسم له مربّع رقم فارغ)
\newcommand{\plainsec}[1]{%
  \par\vspace{6pt}%
  {\large\bfseries\color{primary} #1}\par
  \vspace{-0.35em}{\color{primary}\hrule height 1.2pt}\par\vspace{3pt}%
}

% ── Reusable boxes ───────────────────────────────────────────────────────────
% صندوق معلومات عام
\newtcolorbox{infobox}[1][]{
  enhanced, colback=lightbg, colframe=secondary,
  coltitle=white, fonttitle=\bfseries, title=#1,
  boxrule=0.9pt, arc=4pt, left=10pt, right=10pt, top=8pt, bottom=8pt,
  breakable,
}
% صندوق "أتذكّر" — الخلاصة
\newtcolorbox{rememberbox}[1][أتذكّر]{
  enhanced, colback=lightbg, colframe=primary,
  coltitle=white, fonttitle=\bfseries, title=#1,
  boxrule=1pt, arc=4pt, left=10pt, right=10pt, top=8pt, bottom=8pt,
  breakable,
}
% صندوق نشاط
\newtcolorbox{activitybox}[1][نشاط]{
  enhanced, colback=white, colframe=secondary,
  coltitle=white, fonttitle=\bfseries, title=#1,
  boxrule=1pt, arc=4pt, left=10pt, right=10pt, top=8pt, bottom=8pt,
  breakable,
}
% صندوق الإشكالية
\newtcolorbox{problembox}[1][الإشكالية]{
  enhanced, colback=white, colframe=primary,
  coltitle=white, fonttitle=\bfseries, title=#1,
  boxrule=1pt, arc=4pt, left=10pt, right=10pt, top=8pt, bottom=8pt,
  breakable,
}
% صندوق الواجب المنزلي
\newtcolorbox{homeworkbox}[1][واجب منزلي]{
  enhanced, colback=lightbg, colframe=success,
  coltitle=white, fonttitle=\bfseries, title=#1,
  boxrule=1pt, arc=4pt, left=10pt, right=10pt, top=8pt, bottom=8pt,
  breakable,
}

% ── Lists ────────────────────────────────────────────────────────────────────
\newlist{numlist}{enumerate}{1}
\setlist[numlist,1]{
  label=\arabic*.\hspace{0.8em},
  labelwidth=1.4em, left=0em .. 1.4em,
  itemsep=4pt, parsep=0pt, itemindent=0pt,
}
\newlist{bullets}{itemize}{1}
\setlist[bullets,1]{
  label=\textcolor{secondary}{\textbullet},
  left=0.2em .. 1.2em,
  itemsep=3pt, parsep=0pt,
}

% ── Helpers ──────────────────────────────────────────────────────────────────
% سطر نقاط للإجابة (يجب أن يكون داخل صندوق ليرتسم)
\newcommand{\answerline}{\par\vspace{1pt}\noindent\makebox[\linewidth]{\dotfill}}
% عدّة أسطر إجابة
\newcommand{\answerlines}[1]{\begingroup\count255=0 \loop\ifnum\count255<#1 \answerline\advance\count255 by 1 \repeat\endgroup}
% فراغ داخل السطر
\newcommand{\blank}[1][2.2cm]{\underline{\hspace{#1}}}
% مربّع اختيار
\newcommand{\cb}{\raisebox{0.15ex}{\fbox{\rule{0pt}{1.1ex}\hspace{1em}}}}

% ── Lesson title block ───────────────────────────────────────────────────────
% \lessonheader{نوع الوثيقة}{عنوان الدرس}{المجال التعلمي}{الوحدة التعليمية}
\newcommand{\lessonheader}[4]{%
  \begin{center}
    {\small\color{secondary} الجمهورية الجزائرية الديمقراطية الشعبية}\\[1pt]
    {\small\color{secondary} وزارة التربية الوطنية \quad$\cdot$\quad مديرية التربية لولاية أدرار \quad$\cdot$\quad ثانوية الشهيد حكومي العيد أدرار}\\[5pt]
    {\Large\bfseries\color{primary} #1}\\[3pt]
    {\large\bfseries\color{primary} #2}
  \end{center}
  \vspace{-2pt}
  {\color{secondary}\hrule height 1.2pt}
  \vspace{4pt}
  \begin{center}
  \begin{tabular}{@{}>{\bfseries\color{primary}}r@{\hspace{0.5em}}l@{\hspace{2em}}>{\bfseries\color{primary}}r@{\hspace{0.5em}}l@{}}
    المجال التعلمي : & #3 & الوحدة التعليمية : & #4 \\
    المستوى : & السنة الأولى ثانوي & المدة : & ساعة واحدة \\
  \end{tabular}
  \end{center}
  \vspace{2pt}
  {\color{secondary}\hrule height 1.2pt}
  \vspace{6pt}
}

\hypersetup{
  colorlinks=false,
  pdfauthor={لفقير عبدالجليل},
}
